% ============================================================
%  CHAPTER 4 — MACHINE LEARNING PIPELINE
% ============================================================
\chapter{Machine Learning Pipeline}

\section{Overview}

The Aegis-IoT ML pipeline follows a three-phase architecture: (1)~offline training in Python using scikit-learn, (2)~Abstract Syntax Tree transpilation to native Java using \texttt{m2cgen}, and (3)~real-time inference execution within the AnyLogic JVM. This chapter details each phase.

\section{Data Preprocessing}

\subsection{Source Data}

The training data originates from the RT-IoT2022 dataset, preprocessed into \texttt{clean\_mqtt\_traffic.csv} containing 4,132 MQTT flow records. Each record contains:

\begin{lstlisting}[language=Python, caption={CSV Schema}]
"id","packetSize","interArrival","flowDuration","arrivalTimestamp"
"0","7.714286","2462430.605522","32.011598","2026-04-25 00:00:02"
\end{lstlisting}

\subsection{Feature Engineering}

Two feature sets are constructed for different model targets:

\begin{itemize}
    \item \textbf{Latency Prediction Features ($\mathbf{X}_{rf}$):} $[\text{packetSize}, \text{interArrival}]$ — 2-dimensional input to the Random Forest regressor.
    \item \textbf{Anomaly Detection Features ($\mathbf{X}_{svm}$):} $[\text{packetSize}, \text{interArrival}, \text{flowDuration}]$ — 3-dimensional input to the OneClassSVM, requiring the full feature space to characterize normal traffic boundaries.
    \item \textbf{Forecasting Features ($\mathbf{X}_{ts}$):} Sliding windows of 3 consecutive \texttt{flowDuration} values, predicting the value at $t+10$.
\end{itemize}

\section{Model 1: Random Forest Latency Predictor}

\subsection{Architecture}

The Random Forest Regressor is configured with the following hyperparameters:

\begin{table}[H]
\centering
\caption{Random Forest Hyperparameters}
\label{tab:rf_params}
\begin{tabular}{@{}ll@{}}
\toprule
\textbf{Parameter} & \textbf{Value} \\
\midrule
\texttt{n\_estimators} & 10 \\
\texttt{max\_depth} & 10 \\
\texttt{random\_state} & 42 \\
\texttt{test\_size} & 0.2 \\
\bottomrule
\end{tabular}
\end{table}

\subsection{Training Process}

\begin{lstlisting}[language=Python, caption={Random Forest Training (train\_both.py)}]
from sklearn.ensemble import RandomForestRegressor
from sklearn.model_selection import train_test_split

X_lat = df[["packetSize", "interArrival"]].copy()
y_lat = df["flowDuration"].copy()

X_train, X_test, y_train, y_test = train_test_split(
    X_lat, y_lat, test_size=0.2, random_state=42
)

rf = RandomForestRegressor(n_estimators=10, max_depth=10,
                           random_state=42)
rf.fit(X_train, y_train)
\end{lstlisting}

\subsection{Evaluation}

The model is evaluated using Mean Absolute Error (MAE) on the held-out test set:

\begin{equation}
\text{MAE} = \frac{1}{n}\sum_{i=1}^{n}|y_i - \hat{y}_i|
\end{equation}

\section{Model 2: OneClassSVM Anomaly Detector}

\subsection{Architecture}

The OneClassSVM is configured with $\nu = 0.05$, establishing a decision boundary that expects approximately 5\% of the training data to fall outside the normal region:

\begin{lstlisting}[language=Python, caption={OneClassSVM Training}]
from sklearn.svm import OneClassSVM

X_anom = df[["packetSize", "interArrival",
             "flowDuration"]].copy()
ocsvm = OneClassSVM(nu=0.05)
ocsvm.fit(X_anom)
\end{lstlisting}

\subsection{Decision Function}

The OneClassSVM's \texttt{score} method returns a signed distance from the decision boundary. Negative scores indicate anomalous observations:

\begin{equation}
f(x) = \text{sgn}\left(\sum_{i} \alpha_i K(x_i, x) - \rho\right)
\end{equation}

In the Aegis-IoT pipeline, any packet with $\text{anomalyScore} < 0$ is flagged as a potential DDoS intrusion and visualized as a red marker on the Security Radar.

\section{Model 3: Autoregressive Time-Series Forecaster}

\subsection{Sliding Window Architecture}

The forecaster uses a sliding window of 3 consecutive latency measurements to predict the latency 10 time steps into the future:

\begin{lstlisting}[language=Python, caption={Forecaster Training (train\_forecaster.py)}]
latencies = df['flowDuration'].values
X, y = [], []

for i in range(len(latencies) - 13):
    window = latencies[i:i+3]
    target = latencies[i+13]  # 10 steps ahead
    X.append(window)
    y.append(target)

rf_forecast = RandomForestRegressor(
    n_estimators=10, max_depth=10, random_state=42
)
rf_forecast.fit(X_train, y_train)
\end{lstlisting}

\subsection{Real-Time Forecasting}

During simulation execution, the forecaster is invoked when the latency history buffer contains at least 3 entries. The last 3 actual latencies form the input window:

\begin{lstlisting}[language=Java, caption={Forecasting Hook in Cloud\_Received}]
if (latDash.actualLatencies.size() >= 3) {
    int s = latDash.actualLatencies.size();
    double[] window = new double[]{
        latDash.actualLatencies.get(s-3),
        latDash.actualLatencies.get(s-2),
        latDash.actualLatencies.get(s-1)
    };
    double futureLatency = FutureForecaster.score(window);
    latDash.addForecast(futureLatency);
}
\end{lstlisting}

\section{m2cgen: AST Transpilation to Java}

\subsection{The Zero-Latency Inference Problem}

Traditional ML deployment strategies involve serializing trained models (e.g., using \texttt{joblib} or \texttt{pickle}), loading them into a Python runtime, and exposing predictions via REST APIs or socket-based inter-process communication (IPC). This architecture introduces three sources of latency:

\begin{enumerate}
    \item \textbf{Serialization/Deserialization:} Converting model parameters between binary formats adds 1--10ms per call.
    \item \textbf{IPC Overhead:} Socket communication between the JVM and Python runtime adds 5--50ms per prediction.
    \item \textbf{GIL Contention:} Python's Global Interpreter Lock prevents true parallel inference under multi-threaded workloads.
\end{enumerate}

\subsection{The Transpilation Solution}

The \texttt{m2cgen} framework eliminates all three latency sources by transpiling the model's Abstract Syntax Tree (AST) directly into native Java source code:

\begin{lstlisting}[language=Python, caption={m2cgen Export Process}]
import m2cgen as m2c

# Export Random Forest to Java
rf_java = m2c.export_to_java(rf,
    class_name="OfflineAiPredictor")
with open("OfflineAiPredictor.java", "w") as f:
    f.write("package finalccnproject;\n\n" + rf_java)

# Export OneClassSVM to Java
ocsvm_java = m2c.export_to_java(ocsvm,
    class_name="AnomalyModel")
with open("AnomalyModel.java", "w") as f:
    f.write("package finalccnproject;\n\n" + ocsvm_java)
\end{lstlisting}

\subsection{Generated Java Code Structure}

The transpiled \texttt{OfflineAiPredictor.java} file contains 19,216 lines of pure Java code. Each of the 10 Random Forest estimators is represented as a deeply nested \texttt{if-else} tree. The final prediction is computed as:

\begin{equation}
\hat{y} = \frac{1}{T}\sum_{t=1}^{T} h_t(\mathbf{x})
\end{equation}

where $T=10$ is the number of trees and $h_t(\mathbf{x})$ is the prediction of tree $t$. A representative excerpt:

\begin{lstlisting}[language=Java, caption={Transpiled Decision Tree (excerpt)}]
public class OfflineAiPredictor {
    public static double score(double[] input) {
        double var0;
        if (input[1] <= 2928742.375) {
            if (input[1] <= 1466804.5625) {
                if (input[1] <= 854061.3125) {
                    var0 = 2.141209;
                } else {
                    var0 = 3.936286;
                }
            } // ... 19,000+ lines of decision logic
        }
        // Average across all 10 trees
        return (var0 + var1 + ... + var9) / 10.0;
    }
}
\end{lstlisting}

This pure Java implementation executes as native JVM bytecode with zero external dependencies, achieving sub-millisecond inference latency.

\subsection{Voting Regressor (Alternative Pipeline)}

An alternative training pipeline (\texttt{train\_latency\_voting\_regressor.py}) implements a Voting Regressor ensemble combining three heterogeneous base models:

\begin{enumerate}
    \item \textbf{Linear Regression:} Captures linear relationships between features and latency.
    \item \textbf{SVR (Support Vector Regression):} Models non-linear relationships using RBF kernels, wrapped in a \texttt{StandardScaler} pipeline.
    \item \textbf{AdaBoost with Decision Trees:} Boosted ensemble of depth-4 trees for capturing complex feature interactions.
\end{enumerate}

\begin{lstlisting}[language=Python, caption={Voting Regressor Configuration}]
voting_regressor = VotingRegressor(
    estimators=[
        ("linear_regression", LinearRegression()),
        ("scaled_svr", make_pipeline(
            StandardScaler(), SVR())),
        ("adaboost_tree", AdaBoostRegressor(
            estimator=DecisionTreeRegressor(max_depth=4),
            n_estimators=50))
    ]
)
\end{lstlisting}

\section{GAN-Based DDoS Traffic Generation}

To validate the anomaly detection pipeline, a synthetic DDoS traffic generator (\texttt{generate\_ddos.py}) creates 500 adversarial packets:

\begin{lstlisting}[language=Python, caption={GAN-Inspired DDoS Generator}]
# Volumetric DDoS: massive packets, micro-burst timing
sizes = np.random.normal(loc=15000, scale=2000, size=500)
arrivals = np.random.exponential(scale=10.0, size=500)
durations = np.random.normal(loc=150.0, scale=30.0, size=500)
\end{lstlisting}

During simulation runtime, the GAN injection operates inline within the \texttt{Cloud\_Received} hook with a 5\% probability:

\begin{lstlisting}[language=Java, caption={Runtime GAN Injection}]
if (Math.random() < 0.05) {
    agent.packet_size = 15000 + Math.random() * 5000;
    agent.inter_arrival = Math.random() * 5;
    agent.flow_duration = 150.0 + Math.random() * 50;
}
\end{lstlisting}
